\section{Formal Verification and Reproducibility}
\label{sec:formalization}
The formal theorem is
\begin{center}
\texttt{NSP14.NumericalSemigroup.}\\
\texttt{p14\_totalBetti\_le\_intervalCompletion}.
\end{center}
Its parameters are a numerical semigroup $S$, an arbitrary field $k$,
and a natural number $i$. The source and target are the associated
graded algebras of the actual supported formal-power-series semigroup
rings. The presentation variables are their respective minimal atoms.
The invariant \texttt{totalBetti} is defined as the dimension of the
full-variable Koszul homology, and a proved isomorphism identifies this
homology with ordinary $\Tor$ over that polynomial ring.

The source used for this manuscript is the repository
\href{https://github.com/stksk310/NS-P14-Lean}{\texttt{stksk310/NS-P14-Lean}},
at commit
\begin{center}\small
\texttt{a5e122a4d09536b1609c5faa8d419f8faf81671f}.
\end{center}
The compiler is Lean \texttt{v4.36.0-rc2}. The mathlib revision is
\begin{center}\small
\texttt{139105d28c7685a154cb3093f79edb2e16708099}.
\end{center}
The principal interfaces are organized as follows.
\begin{center}\small
\begin{tabular}{p{.28\linewidth}p{.62\linewidth}}
\toprule
Mathematical role&Lean module\\\midrule
Completed-ring bridge&\raggedright\arraybackslash
\texttt{Semigroup.FormalTangentAlgebra},
\texttt{Homological.ActualCoefficientBaseChange}\\
Exact target&\texttt{ArithmeticSequence.BettiFormula},
\texttt{Main.IntervalTarget}\\
Ordinary comparison&\texttt{Ordinary.Closure}\\
Endpoint mechanism&\texttt{Endpoint} and its homological consumers\\
Narrow-width comparison&\texttt{Width.Closure}, \texttt{NarrowFour.Comparison}\\
Small generators&\texttt{SmallEmbedding.Base},
\texttt{SmallEmbedding.ThreeComparison}\\
Exhaustion and assembly&\texttt{Main.Partition}, \texttt{Main.Branches},
\texttt{Main}\\\bottomrule
\end{tabular}
\end{center}
All module names have prefix \texttt{NSP14}.
The released completion record is
\path{verification/run20/RUN20_COMPLETE_20261116_R03.json}; it binds
the proof snapshot, validation evidence, and a separate AI-assisted audit.
The recorded transitive axioms are Lean's standard
\texttt{propext}, \texttt{Classical.choice}, and \texttt{Quot.sound}.
Published mathematical inputs used in the proof of
Theorem~\ref{thm:main} are proved, in the forms required there, in the
formal development. Corollary~\ref{cor:width-bound} additionally uses
the standard filtered-Koszul comparison with the complete local ring.
The repository's original fifty-entry dependency index and its final
paper-correspondence ledger give finer declaration-level references.

Lean kernel checking is distinct from the generative-AI audit workflow:
the formal theorem is checked against its definitions and proof terms.
The accompanying Python scripts orchestrate compilation, declaration
and axiom checks, and the preservation of verification evidence.
Finite arithmetic checks supplement the general proofs and are not
used to extrapolate an unbounded assertion from sampled cases.

To reproduce the released verification, retain the repository's Git
history and check out the proof snapshot identified by the completion
record, rather than the later commit that packages the release:
\begin{center}\small
\texttt{1a5c6bac9a4fd38a8bc5d13b4c7372e654466e6e}.
\end{center}
Install the pinned compiler and dependencies, including their build
caches. With the pinned Lean executables on \texttt{PATH}, run
\begin{small}
\begin{verbatim}
python3 scripts/verify_run20.py \
  --fresh-project-cache \
  --output-dir verification/run20/runs/independent_new_id
\end{verbatim}
\end{small}
Choose a new output directory for each run. The verifier rebuilds the
project library, intermediate code, tests, and audits while retaining
the fixed dependency cache. The audit archive does not contain an
offline compiler and dependency installation. The archived release
evidence and a new execution of these instructions are distinct records.
